\section{Positive representing densities}\label{sec:densities}

We now construct the representing probabilities required by
Lemma~\ref{peaks:mean-criterion}.  Each step adds an oscillation on
the current patches.  A radial integration by parts makes every
holomorphic moment of that oscillation vanish exactly.  The
separation between consecutive degrees will ensure that the new
density stays positive.

Keep the profiles $f,b,R_0$ from Lemma~\ref{peaks:profiles} fixed,
and put
\[
 c_0=\frac{1-\|f\|_\infty}{4}>0.
\]
Fix any $D_0>\max\{1,2\sqrt{2R_0}\}$.  For an integer $Q\ge2$
large enough that every degree below meets the patch thresholds,
put
\begin{equation}\label{density:degrees}
 k_j=Q^j,\qquad \ell_0=0,\qquad
 \ell_j=\sum_{l=1}^j k_l\quad(j\ge1).
\end{equation}
Thus $\ell_{j-1}<k_j$ and $\ell_j<2k_j$.  At each degree use any
of the separated sets and patches constructed in Section~\ref{sec:peaks}.

Start with $W_0=1$.  At stage $j$, set $k=k_j$.  For each
$p\in\mathcal P_k$ let
\[
 S_p(\theta)=W_{j-1}(e^{i\theta}p).
\]
To determine the correction, consider a phase mode $e^{is\theta}$
with $|s|\le\ell_{j-1}$, the frequency range maintained below.
Multiplication by $e^{ik\theta}$ shifts it to frequency
$m=k+s>0$; its conjugate can pair with the holomorphic monomial
$\xi_1^m=e^{-my/k}e^{im\theta}$.  After canceling the surface-measure factor
$e^{-2y/k}$, its radial moment will vanish if its coefficient is
$b'-(m/k)b$, because
\[
 e^{-my/k}\left(b'-\frac{m}{k}b\right)
 =\frac{d}{dy}\left(e^{-my/k}b\right),
 \qquad b(0)=b(R_0)=0.
\]
The identity $f=b'-b$ therefore prescribes the coefficient
$f-(s/k)b$ for this mode.  Summing the modes leads to the following
definition in the adapted coordinates of $\mathcal U_{k,p}$:
\begin{equation}\label{density:recursion}
 \begin{split}
 A_{j,p}(\xi)
 &=e^{2y/k}\operatorname{Re}\left[
     e^{ik\theta}\left(f(y)S_p(\theta)
               -b(y)\frac{S_p'(\theta)}{ik}\right)\right],\\
 W_j&=W_{j-1}+\sum_{p\in\mathcal P_k}A_{j,p},
 \end{split}
\end{equation}
where each correction is extended by zero off its patch.  The
prime on $S_p$ denotes differentiation in $\theta$.  This recursion
is defined without assuming positivity; that property will be
proved for one fixed sufficiently large $Q$.

\begin{lemma}[Smoothness of the corrections]\label{density:smooth}
Every $W_j$ defined by \eqref{density:recursion} is a smooth real
function on $S$.
\end{lemma}

\begin{proof}
Proceed by induction, beginning with $W_0=1$.  The restriction
$S_p$ is a smooth periodic function.  To check the central circle
of a patch, use the smooth coordinates
\[
 (\theta,\zeta)\longmapsto
       (e^{i\theta}\sqrt{1-|\zeta|^2},\zeta),
 \qquad y=-\frac{k}{2}\log(1-|\zeta|^2).
\]
All terms of \eqref{density:recursion} are smooth in these
coordinates, including at $\zeta=0$.  They do not depend on the
transverse polar angle $\gamma$, which would be singular there.
Their dependence on $\theta$ is periodic, so it defines a smooth
function along the full central circle.  At the outer boundary,
$f$ and $b$ vanish in a collar of $R_0$; hence the correction is
identically zero in a collar of the patch boundary.  Its extension
by zero is therefore smooth.  Summing the finitely many corrections
proves the induction.
\end{proof}

Let $T$ denote the phase derivative on $S$,
\[
 TW(\xi)=\left.\frac{d}{d\omega}W(e^{i\omega}\xi)
                  \right|_{\omega=0}.
\]
A horizontal vector $Z\in T_\xi S$ will always mean a
\emph{Euclidean unit} real tangent vector satisfying the complex
equation $\langle Z,\xi\rangle=0$.

A function has phase bandwidth at most $\ell$ if, on every orbit
$\omega\mapsto e^{i\omega}\xi$, it is a trigonometric polynomial
with frequencies in $[-\ell,\ell]$.  Its orbitwise mean is the
constant coefficient of this polynomial.

\begin{theorem}[Uniform representing densities]\label{density:main}
There is a constant $K_*>0$, depending only on the fixed profiles
$f,b$ and $R_0$, with the following property.  For every fixed
$D_0>\max\{1,2\sqrt{2R_0}\}$ there is an integer threshold $Q_0$
such that, for every integer $Q\ge Q_0$ and every choice of the
separated sets above, the recursion \eqref{density:recursion}
has the following properties for all $j\ge1$:
\begin{enumerate}
\item $W_j$ is smooth and strictly positive, and $W_j\,d\sigma$
represents holomorphic evaluation at zero.  In particular it is
a probability measure.
\item Under global phase rotation, $W_j$ has frequencies only
between $-\ell_j$ and $\ell_j$, and its mean on every phase orbit
is exactly one.
\item At every point of $S$,
\begin{equation}\label{density:uniform-bounds}
 \begin{gathered}
 c_0W_{j-1}\le W_j\le2W_{j-1},\\
 |ZW_j|\le K_*\sqrt{k_j}\,W_j,
 \qquad |TW_j|\le K_*k_jW_j
 \end{gathered}
\end{equation}
for every unit horizontal vector $Z$.
\end{enumerate}
The threshold $Q_0$ may also be enlarged to impose any fixed lower
threshold on the degrees $k_j$.
\end{theorem}

We first prove the exact moment and frequency assertions; they
hold before positivity is known.  The remainder of the theorem,
including the order of choices of $K_*$ and $Q_0$, is proved in
the next subsection.

\subsection{Exact moments and phase frequencies}

\begin{lemma}[Moment preservation]\label{density:moments}
For every stage of \eqref{density:recursion}, the phase bandwidth
of $W_j$ is at most $\ell_j$, its orbitwise mean is one, and
\begin{equation}\label{density:representing}
 \int_S p(\xi)W_j(\xi)\,d\sigma(\xi)=p(0)
 \quad\text{for every holomorphic polynomial }p.
\end{equation}
Each individual correction has zero integral against every such
polynomial, including the constant polynomial.
\end{lemma}

\begin{proof}
At stage zero the phase assertions are immediate.  The polynomial
identity for $\sigma$ follows by global phase averaging: every
nonconstant homogeneous holomorphic polynomial has nonzero
positive phase frequency.

Assume the assertions through stage $j-1$, and write $k=k_j$.
The central restriction has a finite expansion
\[
 S_p(\theta)=\sum_{|s|\le\ell_{j-1}}\widehat S_p(s) e^{is\theta},
 \qquad \widehat S_p(0)=1.
\]
Since $b'-b=f$, the bracket in \eqref{density:recursion} is
\[
 \sum_s \widehat S_p(s)\left(b'-\frac{k+s}{k}b\right)e^{is\theta}.
\]
All integers $m=k+s$ are positive because $\ell_{j-1}<k$.
The patch and $y$ are invariant under global phase rotation,
whereas $\theta$ increases by the rotation angle.  Hence the
correction has frequencies in
\begin{equation}\label{density:bandwidth}
 [k-\ell_{j-1},k+\ell_{j-1}]
 \ \cup\ [-k-\ell_{j-1},-k+\ell_{j-1}].
\end{equation}
They omit zero.  Multiplying by a function of $y$ and extending
by zero across the invariant patch boundary do not alter these
frequencies.  The bandwidth of $W_j$ is thus at most
$k+\ell_{j-1}=\ell_j$, and its orbitwise mean remains one.

For the full moment calculation, use an adapted monomial
$\xi_1^n\xi_2^l$, with nonnegative integers $n,l$.  If $l>0$,
averaging in the transverse phase $\gamma$ gives zero.  If
$l=n=0$, averaging in $\theta$ gives zero because all $m=k+s$
are positive.  Suppose then that $l=0$ and $n\ge1$.  As $b$ is
real, the correction is
\[
 A_{j,p}=\frac{e^{2y/k}}2
 \sum_s\left(\widehat S_p(s) e^{i(k+s)\theta}
              +\overline{\widehat S_p(s)}e^{-i(k+s)\theta}\right)
       \left(b'-\frac{k+s}{k}b\right).
\]
Since $\xi_1^n=e^{-ny/k}e^{in\theta}$, only the conjugate
frequency with $k+s=n$ can survive the $\theta$ integration.
Formula~\eqref{peaks:patch-measure} gives exactly
\[
 \begin{split}
 \int_{\mathcal U_{k,p}}\xi_1^n A_{j,p}\,d\sigma
 &=\frac{\overline{\widehat S_p(n-k)}}{k}
   \int_0^{R_0}e^{-ny/k}\left(b'-\frac nk b\right)dy\\
 &=\frac{\overline{\widehat S_p(n-k)}}{k}
   \left[e^{-ny/k}b(y)\right]_0^{R_0}=0.
 \end{split}
\]
Here a coefficient outside the previous bandwidth is zero.  The
factor $2/k$ in surface measure combines with the $1/2$ from the
real part to leave $1/k$.  Both endpoint terms vanish because
$b(0)=b(R_0)=0$.
The argument applies to every degree $n$, without a degree cutoff.
Unitary changes of coordinates preserve holomorphic polynomials,
so each correction annihilates every polynomial in the original
coordinates.  Adding the corrections proves
\eqref{density:representing} and completes the induction.
\end{proof}

In particular, every $W_j$ has total integral one.  What remains
for Theorem~\ref{density:main} is to prove that these exact
representing densities are positive with the uniform derivative
bounds.  The next estimates accomplish this with a single choice
of $Q$, independent of the number of stages.
